\setgebiet{4}{Analysis}{tief}
\chapter{Analysis}
\marginnote{\small\itshape\color{grau}\textbf{Sie brauchen:} Band 3, Kapitel 1--13. Aus Band 0 die Funktionen und die Ableitungsregeln (Kap.\ 7--12); aus Band 1 die Induktion.}[12pt]
\worum{Analysis ist die Kunst, mit dem Unendlichen so zu rechnen, dass am Ende eine Zahl steht. Grenzwert, Stetigkeit, Ableitung, Integral --- vier Begriffe, ein Mechanismus: $\varepsilon$, dann $\delta$ oder $N$. Die Aufgaben üben das Rechnen \emph{und} das Begründen, denn in der Analysis ist der zweite Teil der, an dem man scheitert.}

\section{Was Sie hier üben}
\begin{center}
{\footnotesize\sffamily\setlength{\tabcolsep}{4pt}
\renewcommand{\arraystretch}{1.5}
\begin{tabular}{@{}lP{4.9cm}lcll@{}}
\toprule
\textbf{Nr.} & \textbf{Titel} & \textbf{Sorte} & \textbf{Grad} & \textbf{Zeit} & \textbf{Band/Kap.}\\
\midrule
3.1 & Supremum ist nicht Maximum & \typtypisch & \grad1 & 10 min & 3/1\\
3.2 & Supremum und Infimum einer Folge & \typwerk & \grad1 & 5 min & 3/1\\
3.3 & Eine Betragsungleichung & \typwerk & \grad1 & 10 min & 3/1\\
3.4 & Grenzwerte, die man sehen muss & \typtypisch & \grad1 & 10 min & 3/2\\
3.5 & $0\cdot\infty$ und $1^{\infty}$ & \typstolper & \grad2 & 10 min & 3/2\\
3.6 & Ein $\varepsilon$-$N$-Beweis & \typwerk & \grad1 & 10 min & 3/2\\
3.7 & $\sqrt[n]{n}\to1$ & \typwerk & \grad2 & 15 min & 3/2\\
3.8 & Eine rekursive Folge & \typwerk & \grad2 & 15 min & 3/2\\
3.9 & Die harmonische Reihe ist keine Cauchy-Folge & \typwerk & \grad2 & 10 min & 3/2\\
3.10 & Vier Reihen, vier Kriterien & \typtypisch & \grad2 & 15 min & 3/3\\
3.11 & Umordnen verboten & \typstolper & \grad2 & 15 min & 3/3\\
3.12 & Eine Teleskopreihe & \typwerk & \grad1 & 10 min & 3/3\\
3.13 & Die $p$-Reihe & \typwerk & \grad2 & 10 min & 3/3\\
3.14 & Fakultät gegen Potenz & \typwerk & \grad2 & 10 min & 3/3\\
3.15 & Drei Konvergenzradien & \typwerk & \grad2 & 15 min & 3/3\\
3.16 & Stetig machen, Nullstelle finden & \typtypisch & \grad1 & 10 min & 3/4\\
3.17 & Stetig, also beschränkt? & \typstolper & \grad1 & 10 min & 3/4\\
3.18 & $x\sin\frac1x$ & \typwerk & \grad2 & 10 min & 3/4\\
3.19 & Ungerader Grad, reelle Nullstelle & \typwerk & \grad2 & 10 min & 3/4\\
3.20 & Die Wurzel ist gleichmäßig stetig & \typwerk & \grad2 & 15 min & 3/4\\
3.21 & Gleichungen mit $e$ und $\ln$ & \typtypisch & \grad1 & 10 min & 3/5\\
3.22 & $\cos(0{,}5)$ auf vier Stellen & \typwerk & \grad1 & 10 min & 3/5\\
3.23 & Ableiten im Akkord & \typtypisch & \grad1 & 10 min & 3/6\\
3.24 & Stetig, aber nicht differenzierbar & \typstolper & \grad1 & 10 min & 3/6\\
3.25 & Alle drei Regeln auf einmal & \typwerk & \grad1 & 10 min & 3/6\\
3.26 & Differenzierbar mit unstetiger Ableitung & \typwerk & \grad2 & 15 min & 3/6\\
3.27 & Ungleichungen mit dem Mittelwertsatz & \typtypisch & \grad2 & 15 min & 3/7\\
3.28 & L'Hospital, falsch angewandt & \typstolper & \grad2 & 10 min & 3/7\\
3.29 & Drei Grenzwerte mit L'Hospital & \typwerk & \grad1 & 10 min & 3/7\\
3.30 & Eine Kurvendiskussion & \typwerk & \grad2 & 20 min & 3/7\\
3.31 & Rechnen mit Taylorreihen & \typtypisch & \grad2 & 15 min & 3/8\\
3.32 & $\sqrt{1{,}1}$ mit Fehlerschranke & \typwerk & \grad1 & 10 min & 3/8\\
3.33 & Die Reihe von $\ln(1 + x)$ & \typwerk & \grad2 & 15 min & 3/8\\
3.34 & Summen, die Integrale sind & \typtypisch & \grad2 & 15 min & 3/9\\
3.35 & $\int_0^1x^{2}\,dx$ von Hand & \typwerk & \grad2 & 15 min & 3/9\\
3.36 & Sechs Integrale, sechs Handgriffe & \typtypisch & \grad2 & 20 min & 3/10\\
3.37 & Drei Integrale, drei Fehler & \typstolper & \grad2 & 15 min & 3/10\\
3.38 & Partielle Integration & \typwerk & \grad1 & 15 min & 3/10\\
3.39 & Substitution & \typwerk & \grad2 & 15 min & 3/10\\
3.40 & Partialbruchzerlegung & \typwerk & \grad1 & 10 min & 3/10\\
3.41 & Der Umfang des Einheitskreises & \typwerk & \grad2 & 15 min & 3/10\\
3.42 & $\int x^{-p}$, zweimal uneigentlich & \typwerk & \grad2 & 10 min & 3/11\\
3.43 & $\int\frac{\sin x}x$ & \typwerk & \grad3 & 20 min & 3/11\\
3.44 & Vergleichen statt rechnen & \typtypisch & \grad2 & 15 min & 3/11\\
3.45 & Punktweise ist nicht gleichmäßig & \typstolper & \grad3 & 20 min & 3/12\\
3.46 & $\sum kx^{k}$ & \typwerk & \grad1 & 10 min & 3/12\\
3.47 & Gradient, Richtung, Tangentialebene & \typtypisch & \grad2 & 15 min & 3/13\\
3.48 & Kritische Punkte klassifizieren & \typwerk & \grad2 & 15 min & 3/13\\
3.49 & Lagrange am Rechteck & \typwerk & \grad1 & 10 min & 3/13\\
3.50 & Das Volumen der Kugel & \typwerk & \grad2 & 10 min & 3/13\\
\bottomrule
\end{tabular}}
\end{center}

\section{Woran Sie merken, dass Sie es können}
\begin{itemize}[leftmargin=1.2em,itemsep=2pt]
\item Sie können zu einem gegebenen $\varepsilon$ ein $N$ oder $\delta$ \emph{ausrechnen} --- nicht nur sagen, dass es eines gibt.
\item Sie wählen zu einer Reihe das passende Kriterium in zehn Sekunden und wissen, dass $a_k\to0$ nichts beweist.
\item Sie kennen die Voraussetzungen der großen Sätze (Maximum, Zwischenwert, Mittelwert, L'Hospital, Vertauschung von Grenzwert und Integral) und ein Gegenbeispiel für jede, die fehlt.
\item Sie integrieren mit den vier Techniken und erkennen ein uneigentliches Integral, bevor Sie einsetzen.
\end{itemize}

\section{Aufgaben}

%%%%%%%%%%% Kapitel 1
\stupser{tief}{Supremum: kleinste obere Schranke. Es gehört zur Menge, wenn es ein Maximum gibt --- und sonst nicht.}{22pt}
\begin{karte}{tief}{3.1}{\typtypisch}{\meta{1}{10 min}{Band 3, Kap.\ 1}}
\kartentitel{Supremum ist nicht Maximum}
\worum{Der erste Begriff der Analysis, und der erste, den man falsch verwendet: Das Supremum muss nicht zur Menge gehören.}
\methoden{Supremum und Infimum, Epsilon-Beweis, Vollständigkeit}
\begin{enumerate}[label=(\alph*),leftmargin=1.6em,itemsep=2pt]
\item Bestimmen Sie Supremum, Infimum, Maximum und Minimum (sofern vorhanden) von $A = \{1 - \frac1n\colon n\ge1\}$, $B = (0, 1)$, $C = [0, 1]$, $D = \{x\in\mathbb Q\colon x^{2} < 2\}$.
\item Zeigen Sie für $A$ mit der $\varepsilon$-Charakterisierung: Zu jedem $\varepsilon > 0$ gibt es ein $a\in A$ mit $a > 1 - \varepsilon$.
\item Warum hat $D$ in $\mathbb Q$ kein Supremum, in $\mathbb R$ aber schon? Welches Axiom steckt dahinter?
\end{enumerate}
\end{karte}

\stupser{gruen}{$\frac{n - 1}{n + 1} = 1 - \frac2{n + 1}$: wächst monoton, startet bei $0$, nähert sich $1$.}{22pt}
\begin{karte}{gruen}{3.2}{\typwerk}{\meta{1}{5 min}{Band 3, Kap.\ 1, Aufgabe 1.1}}
\kartentitel{Supremum und Infimum einer Folge}
\worum{Die Menge der Folgenglieder: Was ist die kleinste obere, was die größte untere Schranke --- und welche wird angenommen?}
\methoden{Supremum und Infimum, Polynomdivision, Monotonie}
Bestimmen Sie Supremum und Infimum von $M = \bigl\{\frac{n - 1}{n + 1}\colon n\in\mathbb N, n\ge1\bigr\}$. Werden sie angenommen?
\end{karte}

\stupser{gruen}{Die Beträge wechseln bei $x = \frac32$ und $x = -1$. Drei Bereiche, drei Rechnungen --- oder quadrieren, weil beide Seiten nichtnegativ sind.}{22pt}
\begin{karte}{gruen}{3.3}{\typwerk}{\meta{1}{10 min}{Band 3, Kap.\ 1, Aufgabe 1.2}}
\kartentitel{Eine Betragsungleichung}
\worum{Fallunterscheidung ist der sichere Weg, Quadrieren der schnelle --- wenn man weiß, wann es erlaubt ist.}
\methoden{Betrag, Fallunterscheidung, Quadrieren}
Lösen Sie $|2x - 3| < |x + 1|$ auf zwei Arten und vergleichen Sie.
\end{karte}

%%%%%%%%%%% Kapitel 2
\stupser{tief}{Höchste Potenz ausklammern; bei Wurzeldifferenzen mit der Summe erweitern; bei $3^{n}$ gegen $2^{n}$ durch $3^{n}$ teilen.}{22pt}
\begin{karte}{tief}{3.4}{\typtypisch}{\meta{1}{10 min}{Band 3, Kap.\ 2}}
\kartentitel{Grenzwerte, die man sehen muss}
\worum{Vier Grundformen, die in jeder Klausur stehen. Wer die Handgriffe kennt, braucht keine $\varepsilon$ mehr --- die Grenzwertsätze haben sie schon erledigt.}
\methoden{Grenzwerte, Ausklammern}
Bestimmen Sie die Grenzwerte, falls sie existieren:
\begin{align*} &\text{(a) } \frac{2n^{2} + 3}{n^{2} - n}, &&\text{(b) } \sqrt{n^{2} + n} - n,\\ &\text{(c) } \frac{3^{n} + 2^{n}}{3^{n} - 2^{n}}, &&\text{(d) } (-1)^{n}\frac n{n + 1}. \end{align*}
\end{karte}

\stupser{akzent}{$n\cdot\frac1n = 1$, $n^{2}\cdot\frac1n = n$, $\frac1n\cdot\frac1n\cdot n = \frac1n$: Dreimal \enquote{$\infty\cdot0$}, drei verschiedene Ergebnisse.}{22pt}
\begin{karte}{akzent}{3.5}{\typstolper}{\meta{2}{10 min}{Band 3, Kap.\ 2}}
\kartentitel{$0\cdot\infty$ und $1^{\infty}$}
\worum{Die Grenzwertsätze gelten für konvergente Folgen. Wer sie auf $\infty$ anwendet, bekommt \enquote{$0\cdot\infty = 0$} und \enquote{$1^{\infty} = 1$} --- und beides ist falsch.}
\methoden{Unbestimmte Formen, Gegenbeispiel, Grenzwerte}
\begin{enumerate}[label=(\alph*),leftmargin=1.6em,itemsep=2pt]
\item Bestimmen Sie $\lim a_nb_n$ für $a_n = \frac1n$ und $b_n = n$, $b_n = n^{2}$, $b_n = \sqrt n$, $b_n = (-1)^{n}n$. Was zeigt das über \enquote{$0\cdot\infty$}?
\item Jemand argumentiert: $\bigl(1 + \frac1n\bigr)^{n}\to1^{\infty} = 1$. Berechnen Sie die Folge für $n = 1, 10, 1000$ und finden Sie den Fehler.
\item Formulieren Sie die Grenzwertsätze für Produkt und Quotient mit ihren \emph{Voraussetzungen}.
\item Welche der folgenden Formen sind bestimmt, welche unbestimmt: $\frac0\infty$, $\frac\infty0$, $\infty - \infty$, $0^{\infty}$, $\infty^{0}$, $\frac{1}{0^{+}}$?
\end{enumerate}
\end{karte}

\stupser{gruen}{$\bigl|\frac{3n + 2}{2n + 5} - \frac32\bigr| = \frac{11}{2(2n + 5)}$. Das soll kleiner als $\varepsilon$ sein --- nach $n$ auflösen.}{22pt}
\begin{karte}{gruen}{3.6}{\typwerk}{\meta{1}{10 min}{Band 3, Kap.\ 2, Aufgabe 2.1}}
\kartentitel{Ein $\varepsilon$-$N$-Beweis}
\worum{Einmal die Definition in voller Länge: Abstand ausrechnen, nach $n$ auflösen, $N$ hinschreiben.}
\methoden{Epsilon-Beweis, Grenzwerte}
Zeigen Sie mit der Definition, dass $\lim_{n\to\infty}\dfrac{3n + 2}{2n + 5} = \dfrac32$. Geben Sie für $\varepsilon = 0{,}01$ ein konkretes $N$ an.
\end{karte}

\stupser{gruen}{Schreiben Sie $\sqrt[n]n = 1 + h_n$ mit $h_n\ge0$. Dann $n = (1 + h_n)^{n}\ge\binom n2h_n^{2}$.}{22pt}
\begin{karte}{gruen}{3.7}{\typwerk}{\meta{2}{15 min}{Band 3, Kap.\ 2, Aufgabe 2.2}}
\kartentitel{$\sqrt[n]{n}\to1$}
\worum{Ein berühmter Grenzwert mit einem berühmten Trick: den Überschuss über $1$ benennen und mit dem binomischen Satz einfangen.}
\methoden{Binomischer Satz, Einschließungssatz, Überschuss benennen}
Bestimmen Sie $\lim_{n\to\infty}\sqrt[n]n$ und beweisen Sie das Ergebnis.
\end{karte}

\stupser{gruen}{Erst zeigen: $a_n < 2$ für alle $n$ (Induktion). Dann: $a_n < 2\Rightarrow a_n^{2} < 2 + a_n\Rightarrow a_n < a_{n+1}$.}{22pt}
\begin{karte}{gruen}{3.8}{\typwerk}{\meta{2}{15 min}{Band 3, Kap.\ 2, Aufgabe 2.3}}
\kartentitel{Eine rekursive Folge}
\worum{Das Standardverfahren für rekursive Folgen: monoton und beschränkt, also konvergent; dann den Grenzwert aus der Fixpunktgleichung.}
\methoden{Rekursive Folge, Monotonie, Fixpunktgleichung}
$a_1 = 1$, $a_{n+1} = \sqrt{2 + a_n}$. Zeigen Sie, dass die Folge konvergiert, und bestimmen Sie den Grenzwert.
\end{karte}

\stupser{gruen}{$a_{2n} - a_n = \frac1{n + 1} + \dots + \frac1{2n}$: $n$ Summanden, jeder mindestens $\frac1{2n}$.}{22pt}
\begin{karte}{gruen}{3.9}{\typwerk}{\meta{2}{10 min}{Band 3, Kap.\ 2, Aufgabe 2.4}}
\kartentitel{Die harmonische Reihe ist keine Cauchy-Folge}
\worum{Benachbarte Glieder rücken zusammen, und trotzdem divergiert die Folge. Das Gegenbeispiel, das man zur Cauchy-Definition kennen muss.}
\methoden{Cauchy-Folge, Harmonische Reihe, Abschätzung nach unten}
$a_n = \sum_{k=1}^{n}\frac1k$. Zeigen Sie, dass $(a_n)$ keine Cauchy-Folge ist (und damit divergiert), obwohl $a_{n+1} - a_n\to0$.
\end{karte}

%%%%%%%%%%% Kapitel 3
\stupser{tief}{Rationale Glieder: Vergleich mit $\frac1{k^{p}}$. Alternierend: Leibniz. Fakultäten und Potenzen: Quotient. Immer zuerst: Geht $a_k$ gegen null?}{22pt}
\begin{karte}{tief}{3.10}{\typtypisch}{\meta{2}{15 min}{Band 3, Kap.\ 3}}
\kartentitel{Vier Reihen, vier Kriterien}
\worum{Reihen untersuchen heißt: das Kriterium wählen, das zur Gestalt der Glieder passt. Hier vier Gestalten.}
\methoden{Konvergenzkriterien, Vergleichskriterium, Leibniz-Kriterium, Quotientenkriterium}
Untersuchen Sie auf Konvergenz und absolute Konvergenz:
\begin{align*} &\text{(a) } \sum_{k\ge1}\frac k{k^{2} + 1}, &&\text{(b) } \sum_{k\ge1}\frac1{k^{2} + 1},\\ &\text{(c) } \sum_{k\ge1}\frac{(-1)^{k}}{\sqrt k}, &&\text{(d) } \sum_{k\ge1}\frac{k^{2}}{2^{k}}. \end{align*}
\end{karte}

\stupser{akzent}{Nehmen Sie die alternierende harmonische Reihe $S = 1 - \frac12 + \frac13 - \dots$ und daneben $\frac S2 = 0 + \frac12 + 0 - \frac14 + 0 + \frac16 - \dots$. Addieren Sie gliedweise.}{22pt}
\begin{karte}{akzent}{3.11}{\typstolper}{\meta{2}{15 min}{Band 3, Kap.\ 3}}
\kartentitel{Umordnen verboten}
\worum{Bei endlichen Summen darf man umsortieren. Bei unendlichen nur, wenn sie absolut konvergieren --- sonst kann man jeden Wert erzeugen. Dazu der ältere Irrtum: \enquote{$a_k\to0$, also konvergiert die Reihe}.}
\methoden{Absolute Konvergenz, Umordnung, Gegenbeispiel}
\begin{enumerate}[label=(\alph*),leftmargin=1.6em,itemsep=2pt]
\item $S = \sum_{k\ge1}\frac{(-1)^{k+1}}k = \ln2$ (Aufgabe 3.33). Zeigen Sie, dass die Umordnung $1 + \frac13 - \frac12 + \frac15 + \frac17 - \frac14 + \dots$ (zwei positive, ein negatives Glied) den Wert $\frac32\ln2$ hat.
\item Warum ist das kein Widerspruch? Welche Eigenschaft fehlt der Reihe?
\item Nennen Sie eine Reihe mit $a_k\to0$, die divergiert, und eine konvergente Reihe, die nicht absolut konvergiert.
\item Was besagt der Riemannsche Umordnungssatz?
\end{enumerate}
\end{karte}

\stupser{gruen}{$\frac1{k(k + 1)} = \frac1k - \frac1{k + 1}$. In der Partialsumme hebt sich fast alles weg.}{22pt}
\begin{karte}{gruen}{3.12}{\typwerk}{\meta{1}{10 min}{Band 3, Kap.\ 3, Aufgabe 3.1}}
\kartentitel{Eine Teleskopreihe}
\worum{Eine der wenigen Reihen, deren Summe man exakt ausrechnen kann --- weil die Partialsummen zusammenklappen.}
\methoden{Teleskopsumme, Partialbruchzerlegung}
Untersuchen Sie $\sum_{k\ge1}\frac1{k(k + 1)}$ auf Konvergenz und bestimmen Sie gegebenenfalls die Summe.
\end{karte}

\stupser{gruen}{Verdichtung: $\sum2^{j}a_{2^{j}}$ konvergiert genau dann, wenn $\sum a_k$ konvergiert. Hier wird daraus eine geometrische Reihe mit Quotient $2^{1 - p}$.}{22pt}
\begin{karte}{gruen}{3.13}{\typwerk}{\meta{2}{10 min}{Band 3, Kap.\ 3, Aufgabe 3.2}}
\kartentitel{Die $p$-Reihe}
\worum{Die Vergleichsreihe für alles andere: $\sum k^{-p}$ konvergiert genau für $p > 1$. Hier wird es bewiesen.}
\methoden{Verdichtungskriterium, Geometrische Reihe}
Für welche $p > 0$ konvergiert $\sum_{k\ge1}\frac1{k^{p}}$?
\end{karte}

\stupser{gruen}{$\frac{a_{k+1}}{a_k} = \bigl(\frac k{k + 1}\bigr)^{k} = \frac1{(1 + 1/k)^{k}}\to\frac1e$.}{22pt}
\begin{karte}{gruen}{3.14}{\typwerk}{\meta{2}{10 min}{Band 3, Kap.\ 3, Aufgabe 3.3}}
\kartentitel{Fakultät gegen Potenz}
\worum{Das Quotientenkriterium in seinem Element --- und der Grenzwert $e$ taucht auf, wo man ihn nicht erwartet.}
\methoden{Quotientenkriterium, Zahl e}
Untersuchen Sie $\sum_{k\ge1}\frac{k!}{k^{k}}$ auf Konvergenz.
\end{karte}

\stupser{gruen}{$R = \lim\bigl|\frac{c_k}{c_{k+1}}\bigr|$. Für (a) ist das $\frac{k + 1}k\to1$; am Rand $x = \pm1$ einsetzen.}{22pt}
\begin{karte}{gruen}{3.15}{\typwerk}{\meta{2}{15 min}{Band 3, Kap.\ 3, Aufgabe 3.4}}
\kartentitel{Drei Konvergenzradien}
\worum{Drei Potenzreihen mit Radius $1$, $\infty$ und $0$ --- und am Rand des Konvergenzkreises passiert das Interessante.}
\methoden{Konvergenzradius, Randverhalten, Leibniz-Kriterium}
Bestimmen Sie Konvergenzradius und Randverhalten von
\[ \text{(a) } \sum_{k\ge1}\frac{x^{k}}k,\qquad\text{(b) } \sum_{k\ge0}\frac{x^{k}}{k!},\qquad\text{(c) } \sum_{k\ge0}k!\,x^{k}. \]
\end{karte}

%%%%%%%%%%% Kapitel 4
\stupser{tief}{$\frac{x^{2} - 1}{x - 1} = x + 1$ für $x\neq1$. Für den Zwischenwertsatz: $g(x) = x - \cos x$ an den Intervallenden auswerten.}{22pt}
\begin{karte}{tief}{3.16}{\typtypisch}{\meta{1}{10 min}{Band 3, Kap.\ 4}}
\kartentitel{Stetig machen, Nullstelle finden}
\worum{Zwei Standardaufgaben zur Stetigkeit: eine Lücke schließen und mit dem Zwischenwertsatz eine Lösung nachweisen, die man nicht ausrechnen kann.}
\methoden{Stetige Fortsetzung, Zwischenwertsatz, Intervallhalbierung}
\begin{enumerate}[label=(\alph*),leftmargin=1.6em,itemsep=2pt]
\item Für welches $c$ ist $f(x) = \frac{x^{2} - 1}{x - 1}$ ($x\neq1$), $f(1) = c$ stetig?
\item Für welche $a, b$ ist $g(x) = x^{2}$ für $x\le1$, $g(x) = ax + b$ für $x > 1$ stetig \emph{und} differenzierbar?
\item Zeigen Sie, dass $x = \cos x$ eine Lösung in $(0, \frac\pi2)$ hat, und grenzen Sie sie auf ein Intervall der Länge $0{,}1$ ein.
\item Warum hat die Signumfunktion keine Lösung von $\operatorname{sgn}(x) = \frac12$, obwohl $-1$ und $1$ angenommen werden?
\end{enumerate}
\end{karte}

\stupser{akzent}{$f(x) = \frac1x$ auf $(0, 1)$ ist stetig. Beschränkt? Und $f(x) = x$ auf $(0, 1)$: Wo wäre das Maximum?}{22pt}
\begin{karte}{akzent}{3.17}{\typstolper}{\meta{1}{10 min}{Band 3, Kap.\ 4}}
\kartentitel{Stetig, also beschränkt?}
\worum{Der Satz vom Maximum gilt auf abgeschlossenen, beschränkten Intervallen. Jedes Wort ist nötig, und für jedes weggelassene gibt es ein Gegenbeispiel.}
\methoden{Satz vom Maximum, Gegenbeispiel, Kompaktheit}
\begin{enumerate}[label=(\alph*),leftmargin=1.6em,itemsep=2pt]
\item Geben Sie eine stetige Funktion auf $(0, 1)$ an, die unbeschränkt ist, und eine, die beschränkt ist, aber kein Maximum hat.
\item Geben Sie eine stetige Funktion auf $[0, \infty)$ an, die beschränkt ist und kein Maximum hat.
\item Geben Sie eine Funktion auf $[0, 1]$ an, die kein Maximum hat. Welche Voraussetzung ist verletzt?
\item Formulieren Sie den Satz vom Maximum vollständig und sagen Sie, welche Voraussetzung in (a), (b), (c) jeweils fehlt.
\end{enumerate}
\end{karte}

\stupser{gruen}{$|x\sin\frac1x|\le|x|$: Das zwingt $f$ gegen $0$. Für $g$: Die Punkte $x_n = \frac1{(2n + \frac12)\pi}$ gehen gegen $0$, aber $g(x_n) = 1$.}{22pt}
\begin{karte}{gruen}{3.18}{\typwerk}{\meta{2}{10 min}{Band 3, Kap.\ 4, Aufgabe 4.1}}
\kartentitel{$x\sin\frac1x$}
\worum{Zwei Funktionen, die in jeder Umgebung der Null unendlich oft schwingen --- die eine ist stetig, die andere nicht. Der Faktor $x$ macht den Unterschied.}
\methoden{Epsilon-Beweis, Einschließungssatz, Folgenkriterium}
Untersuchen Sie $f(x) = x\sin\frac1x$ ($x\neq0$), $f(0) = 0$, auf Stetigkeit in $0$. Vergleichen Sie mit $g(x) = \sin\frac1x$ ($x\neq0$), $g(0) = 0$.
\end{karte}

\stupser{gruen}{Für großes $|x|$ verhält sich $p(x)$ wie $a_nx^{n}$ --- und $x^{n}$ mit ungeradem $n$ hat verschiedene Vorzeichen bei $\pm M$.}{22pt}
\begin{karte}{gruen}{3.19}{\typwerk}{\meta{2}{10 min}{Band 3, Kap.\ 4, Aufgabe 4.2}}
\kartentitel{Ungerader Grad, reelle Nullstelle}
\worum{Der Zwischenwertsatz in seiner klassischen Anwendung: Vorzeichenwechsel plus Stetigkeit ergibt eine Nullstelle.}
\methoden{Zwischenwertsatz, Ausklammern}
Zeigen Sie: Jedes Polynom ungeraden Grades hat mindestens eine reelle Nullstelle.
\end{karte}

\stupser{gruen}{$|\sqrt x - \sqrt y|\le\sqrt{|x - y|}$ --- quadrieren Sie das. Dann $\delta = \varepsilon^{2}$.}{22pt}
\begin{karte}{gruen}{3.20}{\typwerk}{\meta{2}{15 min}{Band 3, Kap.\ 4, Aufgabe 4.3}}
\kartentitel{Die Wurzel ist gleichmäßig stetig}
\worum{Gleichmäßig stetig heißt: ein $\delta$ für alle Stellen. Die Wurzel schafft das auf $[0, \infty)$, obwohl ihre Steigung bei $0$ unendlich wird.}
\methoden{Gleichmäßige Stetigkeit, Abschätzung durch Ungleichung}
Zeigen Sie, dass $f(x) = \sqrt x$ auf $[0, \infty)$ gleichmäßig stetig ist.
\end{karte}

%%%%%%%%%%% Kapitel 5
\stupser{tief}{$e^{2x} = (e^{x})^{2}$: Setzen Sie $u = e^{x}$. Bei $\ln$-Gleichungen am Ende prüfen, ob alle Argumente positiv sind.}{22pt}
\begin{karte}{tief}{3.21}{\typtypisch}{\meta{1}{10 min}{Band 3, Kap.\ 5}}
\kartentitel{Gleichungen mit $e$ und $\ln$}
\worum{Rechenroutine mit den elementaren Funktionen: Substitution bei Exponentialgleichungen, Zusammenfassen bei Logarithmen, und die Definitionsbereichsprobe, die man nie vergessen darf.}
\methoden{Substitution (Gleichung), Logarithmieren, Definitionsbereich prüfen}
\begin{enumerate}[label=(\alph*),leftmargin=1.6em,itemsep=2pt]
\item Lösen Sie $e^{2x} - 3e^{x} + 2 = 0$.
\item Lösen Sie $\ln x + \ln(x - 3) = \ln4$.
\item Lösen Sie $2^{x} = 5\cdot3^{x}$.
\item Leiten Sie aus $\exp(a + b) = \exp a\exp b$ die Regel $\ln(xy) = \ln x + \ln y$ her.
\end{enumerate}
\end{karte}

\stupser{gruen}{$\cos x = 1 - \frac{x^{2}}2 + \frac{x^{4}}{24} - \frac{x^{6}}{720} + \dots$; bei alternierenden Reihen ist der Fehler kleiner als das erste weggelassene Glied.}{22pt}
\begin{karte}{gruen}{3.22}{\typwerk}{\meta{1}{10 min}{Band 3, Kap.\ 5, Aufgabe 5.1}}
\kartentitel{$\cos(0{,}5)$ auf vier Stellen}
\worum{Wie der Taschenrechner es macht --- und wie man weiß, wann man aufhören darf.}
\methoden{Potenzreihe, Fehlerschranke}
Berechnen Sie $\cos(0{,}5)$ auf vier Nachkommastellen und geben Sie eine garantierte Fehlerschranke an.
\end{karte}

%%%%%%%%%%% Kapitel 6
\stupser{tief}{$x^{x} = e^{x\ln x}$. $\ln(\cos x)$: Kettenregel mit äußerer Ableitung $\frac1{\cos x}$.}{22pt}
\begin{karte}{tief}{3.23}{\typtypisch}{\meta{1}{10 min}{Band 3, Kap.\ 6}}
\kartentitel{Ableiten im Akkord}
\worum{Fünf Ableitungen, die zusammen alle Regeln abdecken --- und die Umschreibung $x^{x} = e^{x\ln x}$, die man kennen muss.}
\methoden{Ableitungsregeln, Kettenregel, Logarithmieren}
Leiten Sie ab und vereinfachen Sie:
\[ \text{(a) } x^{x},\quad\text{(b) } \ln(\cos x),\quad\text{(c) } \arctan(2x), \]
\[ \text{(d) } \sin^{2}(3x),\quad\text{(e) } xe^{-x^{2}}. \]
\end{karte}

\stupser{akzent}{Bei $|x|$ in $0$: Differenzenquotient von links $-1$, von rechts $+1$. Bei $\sqrt[3]x$ in $0$: $\frac{h^{1/3}}h = h^{-2/3}\to\infty$.}{22pt}
\begin{karte}{akzent}{3.24}{\typstolper}{\meta{1}{10 min}{Band 3, Kap.\ 6}}
\kartentitel{Stetig, aber nicht differenzierbar}
\worum{Differenzierbar impliziert stetig; die Umkehrung gilt nicht, und zwar auf zwei Arten: Knick und senkrechte Tangente. Dazu der Rechenfehler $(fg)' = f'g'$.}
\methoden{Differenzenquotient, Gegenbeispiel, Einseitige Grenzwerte}
\begin{enumerate}[label=(\alph*),leftmargin=1.6em,itemsep=2pt]
\item Zeigen Sie, dass $|x|$ in $0$ stetig, aber nicht differenzierbar ist.
\item Zeigen Sie dasselbe für $\sqrt[3]x$. Was ist hier anders als bei $|x|$?
\item Beweisen Sie: Ist $f$ in $a$ differenzierbar, so ist $f$ in $a$ stetig.
\item Prüfen Sie an $f(x) = g(x) = x$, dass $(fg)'\neq f'g'$, und schreiben Sie die richtige Regel hin.
\end{enumerate}
\end{karte}

\stupser{gruen}{Quotientenregel außen, Produktregel im Zähler, Kettenregel in beiden Faktoren. Sauber notieren, dann zusammenfassen.}{22pt}
\begin{karte}{gruen}{3.25}{\typwerk}{\meta{1}{10 min}{Band 3, Kap.\ 6, Aufgabe 6.1}}
\kartentitel{Alle drei Regeln auf einmal}
\worum{Ein Ausdruck, der Ketten-, Produkt- und Quotientenregel gleichzeitig verlangt. Übung in Buchführung.}
\methoden{Produktregel, Quotientenregel, Kettenregel}
Bestimmen Sie die Ableitung von $f(x) = \dfrac{e^{x^{2}}\sin(3x)}{1 + x^{2}}$.
\end{karte}

\stupser{gruen}{$\frac{f(h) - f(0)}h = h\sin\frac1h\to0$: differenzierbar mit $f'(0) = 0$. Für $x\neq0$ aber $f'(x) = 2x\sin\frac1x - \cos\frac1x$.}{22pt}
\begin{karte}{gruen}{3.26}{\typwerk}{\meta{2}{15 min}{Band 3, Kap.\ 6, Aufgabe 6.2}}
\kartentitel{Differenzierbar mit unstetiger Ableitung}
\worum{Die Ableitung einer differenzierbaren Funktion muss nicht stetig sein. Dieses Beispiel ist das, mit dem man es belegt.}
\methoden{Differenzenquotient, Folgenkriterium, Unstetige Ableitung}
Untersuchen Sie $f(x) = x^{2}\sin\frac1x$ ($x\neq0$), $f(0) = 0$: Ist $f$ in $0$ differenzierbar? Ist $f'$ in $0$ stetig?
\end{karte}

%%%%%%%%%%% Kapitel 7
\stupser{tief}{Differenz der beiden Seiten als Funktion $g$; $g(0) = 0$; Vorzeichen von $g'$ untersuchen. Für (d): Ableitung ist null, also konstant --- Wert bei $x = 1$.}{22pt}
\begin{karte}{tief}{3.27}{\typtypisch}{\meta{2}{15 min}{Band 3, Kap.\ 7}}
\kartentitel{Ungleichungen mit dem Mittelwertsatz}
\worum{Der Mittelwertsatz in Arbeitskleidung: $g' \ge 0\Rightarrow g$ wächst, $g' = 0\Rightarrow g$ konstant. Damit beweist man Ungleichungen und Identitäten, ohne zu rechnen.}
\methoden{Mittelwertsatz}
Beweisen Sie:
\begin{enumerate}[label=(\alph*),leftmargin=1.6em,itemsep=2pt]
\item $e^{x}\ge1 + x$ für alle $x$,
\item $\sin x\le x$ für $x\ge0$,
\item $\ln(1 + x)\le x$ für $x > -1$,
\item $\arctan x + \arctan\frac1x = \frac\pi2$ für $x > 0$.
\end{enumerate}
\end{karte}

\stupser{akzent}{$\lim_{x\to1}\frac{x^{2} + 1}{x + 1}$ ist kein $\frac00$ --- einfach einsetzen. Und $\frac{x + \sin x}x$ für $x\to\infty$: Was passiert mit $\frac{\sin x}x$?}{22pt}
\begin{karte}{akzent}{3.28}{\typstolper}{\meta{2}{10 min}{Band 3, Kap.\ 7}}
\kartentitel{L'Hospital, falsch angewandt}
\worum{Die Regel hat drei Voraussetzungen: unbestimmte Form, Differenzierbarkeit, und der Grenzwert der Ableitungsquotienten muss \emph{existieren}. Drei Beispiele, in denen je eine fehlt.}
\methoden{Regel von L'Hospital, Voraussetzungen prüfen, Gegenbeispiel}
\begin{enumerate}[label=(\alph*),leftmargin=1.6em,itemsep=2pt]
\item Jemand rechnet $\lim_{x\to1}\frac{x^{2} + 1}{x + 1} = \lim\frac{2x}1 = 2$. Richtig?
\item Jemand rechnet $\lim_{x\to\infty}\frac{x + \sin x}x = \lim\frac{1 + \cos x}1$ und schließt: existiert nicht. Richtig?
\item Jemand beweist $\lim_{x\to0}\frac{\sin x}x = 1$ mit L'Hospital. Was ist daran bedenklich, wenn man $(\sin x)' = \cos x$ erst noch herleiten muss?
\item Formulieren Sie die Regel vollständig, mit allen Voraussetzungen.
\end{enumerate}
\end{karte}

\stupser{gruen}{(a) zweimal L'Hospital. (b) $x\ln x = \frac{\ln x}{1/x}$. (c) $x^{x} = e^{x\ln x}$, dann (b).}{22pt}
\begin{karte}{gruen}{3.29}{\typwerk}{\meta{1}{10 min}{Band 3, Kap.\ 7, Aufgabe 7.1}}
\kartentitel{Drei Grenzwerte mit L'Hospital}
\worum{Die drei Formen $\frac00$, $0\cdot\infty$ und $0^{0}$ --- und wie man die letzten beiden auf die erste zurückführt.}
\methoden{Regel von L'Hospital, Unbestimmte Formen, Logarithmieren}
Berechnen Sie (a) $\lim_{x\to0}\dfrac{e^{x} - 1 - x}{x^{2}}$, (b) $\lim_{x\to0^{+}}x\ln x$, (c) $\lim_{x\to0^{+}}x^{x}$.
\end{karte}

\stupser{gruen}{$f(x) = 1 + \frac1{x^{2} - 1}$: Das sieht man die Asymptote $y = 1$ und die Pole sofort.}{22pt}
\begin{karte}{gruen}{3.30}{\typwerk}{\meta{2}{20 min}{Band 3, Kap.\ 7, Aufgabe 7.2}}
\kartentitel{Eine Kurvendiskussion}
\worum{Das vollständige Programm: Definitionsbereich, Symmetrie, Nullstellen, Pole, Asymptoten, Monotonie, Extrema, Krümmung, Skizze.}
\methoden{Kurvendiskussion, Polynomdivision, Krümmung}
Führen Sie eine vollständige Kurvendiskussion durch für $f(x) = \dfrac{x^{2}}{x^{2} - 1}$.
\end{karte}

%%%%%%%%%%% Kapitel 8
\stupser{tief}{$\cos x = 1 - \frac{x^{2}}2 + \frac{x^{4}}{24} - \dots$ einsetzen: Zähler beginnt mit $\frac{x^{4}}{24}$. Für (b) die beiden Reihen multiplizieren und bis $x^{3}$ sammeln.}{22pt}
\begin{karte}{tief}{3.31}{\typtypisch}{\meta{2}{15 min}{Band 3, Kap.\ 8}}
\kartentitel{Rechnen mit Taylorreihen}
\worum{Taylorreihen sind ein Rechenwerkzeug: Grenzwerte ablesen statt L'Hospital, Produkte ausmultiplizieren, Fehler abschätzen.}
\methoden{Potenzreihe, Koeffizientenvergleich, Fehlerschranke}
\begin{enumerate}[label=(\alph*),leftmargin=1.6em,itemsep=2pt]
\item Bestimmen Sie $\lim_{x\to0}\dfrac{\cos x - 1 + \frac{x^{2}}2}{x^{4}}$ und $\lim_{x\to0}\dfrac{\sin x - x}{x^{3}}$ mit Reihen.
\item Bestimmen Sie das Taylorpolynom dritten Grades von $e^{x}\sin x$ um $0$ durch Multiplikation der Reihen.
\item Wie groß ist der Fehler von $\sin x\approx x - \frac{x^{3}}6$ höchstens, wenn $|x|\le0{,}5$?
\end{enumerate}
\end{karte}

\stupser{gruen}{$f(x) = \sqrt x$ um $x_0 = 1$: $f(1) = 1$, $f'(1) = \frac12$, $f''(1) = -\frac14$. Restglied mit $f'''(\xi) = \frac38\xi^{-5/2}\le\frac38$.}{22pt}
\begin{karte}{gruen}{3.32}{\typwerk}{\meta{1}{10 min}{Band 3, Kap.\ 8, Aufgabe 8.1}}
\kartentitel{$\sqrt{1{,}1}$ mit Fehlerschranke}
\worum{Taylorpolynom zweiten Grades plus Lagrange-Restglied: eine Näherung und ihre Garantie.}
\methoden{Taylorpolynom, Fehlerschranke}
Berechnen Sie $\sqrt{1{,}1}$ mit dem Taylorpolynom zweiten Grades und geben Sie eine Fehlerschranke an.
\end{karte}

\stupser{gruen}{$f^{(k)}(0) = (-1)^{k+1}(k - 1)!$, also $\frac{f^{(k)}(0)}{k!} = \frac{(-1)^{k+1}}k$. Am Rand $x = 1$: Leibniz; $x = -1$: harmonisch.}{22pt}
\begin{karte}{gruen}{3.33}{\typwerk}{\meta{2}{15 min}{Band 3, Kap.\ 8, Aufgabe 8.2}}
\kartentitel{Die Reihe von $\ln(1 + x)$}
\worum{Eine Reihe, die am einen Rand konvergiert und am anderen nicht --- und am konvergenten Rand $\ln2$ liefert.}
\methoden{Taylorreihe, Gliedweise Integration, Randverhalten}
Bestimmen Sie die Taylorreihe von $\ln(1 + x)$ um $0$ und ihren Konvergenzbereich. Berechnen Sie damit $\ln2$ als Reihe.
\end{karte}

%%%%%%%%%%% Kapitel 9
\stupser{tief}{$\sum_{k=1}^{n}\frac1{n + k} = \frac1n\sum_{k=1}^{n}\frac1{1 + k/n}$: eine Riemann-Summe von $\frac1{1 + x}$ auf $[0, 1]$.}{22pt}
\begin{karte}{tief}{3.34}{\typtypisch}{\meta{2}{15 min}{Band 3, Kap.\ 9}}
\kartentitel{Summen, die Integrale sind}
\worum{Die Riemann-Summe rückwärts gelesen: Ein Grenzwert von Summen, den man nicht ausrechnen kann, ist ein Integral, das man ausrechnen kann.}
\methoden{Riemann-Summe}
Bestimmen Sie die Grenzwerte für $n\to\infty$:
\[ \text{(a) } \sum_{k=1}^{n}\frac1{n + k},\qquad\text{(b) } \frac1n\sum_{k=1}^{n}\Bigl(\frac kn\Bigr)^{2},\qquad\text{(c) } \sum_{k=1}^{n}\frac n{n^{2} + k^{2}}. \]
\end{karte}

\stupser{gruen}{Obersumme $\frac1{n^{3}}\sum_{i=1}^{n}i^{2} = \frac{n(n + 1)(2n + 1)}{6n^{3}}$ (Aufgabe 1.1 aus Gebiet 1).}{22pt}
\begin{karte}{gruen}{3.35}{\typwerk}{\meta{2}{15 min}{Band 3, Kap.\ 9, Aufgabe 9.1}}
\kartentitel{$\int_0^1x^{2}\,dx$ von Hand}
\worum{Einmal ohne Hauptsatz: Ober- und Untersummen, die zusammenrücken. So war das Integral definiert, bevor es bequem wurde.}
\methoden{Ober- und Untersummen, Summenformeln}
Berechnen Sie $\int_0^1x^{2}\,dx$ direkt über Ober- und Untersummen.
\end{karte}

%%%%%%%%%%% Kapitel 10
\stupser{tief}{Steht die Ableitung des Inneren als Faktor da, substituieren ($\frac x{1 + x^{2}}$, $x\sqrt{1 + x^{2}}$, $\tan x = \frac{\sin x}{\cos x}$). $\sin^{2}x$: Doppelwinkel. $x^{2}e^{x}$: zweimal partiell.}{22pt}
\begin{karte}{tief}{3.36}{\typtypisch}{\meta{2}{20 min}{Band 3, Kap.\ 10}}
\kartentitel{Sechs Integrale, sechs Handgriffe}
\worum{Integrieren ist Mustererkennung. Sechs Muster, die immer wieder vorkommen.}
\methoden{Substitution (Integral), Partielle Integration, Trigonometrische Identität}
\begin{align*} &\text{(a) } \int\frac x{1 + x^{2}}\,dx, &&\text{(b) } \int\tan x\,dx, &&\text{(c) } \int\frac{dx}{x^{2} + 4},\\ &\text{(d) } \int x\sqrt{1 + x^{2}}\,dx, &&\text{(e) } \int\sin^{2}x\,dx, &&\text{(f) } \int x^{2}e^{x}\,dx. \end{align*}
\end{karte}

\stupser{akzent}{$\frac1{x^{2}} > 0$ überall --- kann das Integral $-2$ sein? Bei (b): Die Grenzen gehören zur Variablen, und die Variable wurde gewechselt.}{22pt}
\begin{karte}{akzent}{3.37}{\typstolper}{\meta{2}{15 min}{Band 3, Kap.\ 10}}
\kartentitel{Drei Integrale, drei Fehler}
\worum{Drei Rechnungen, die flüssig aussehen und falsch sind: eine Polstelle übersehen, Grenzen nicht umgerechnet, ein Produkt auseinandergerissen.}
\methoden{Fehler finden, Uneigentliches Integral, Substitution (Integral)}
Finden Sie den Fehler und korrigieren Sie:
\begin{enumerate}[label=(\alph*),leftmargin=1.6em,itemsep=2pt]
\item $\int_{-1}^{1}\frac{dx}{x^{2}} = \bigl[-\frac1x\bigr]_{-1}^{1} = -1 - 1 = -2$.
\item $\int_0^2 2xe^{x^{2}}\,dx$ mit $u = x^{2}$: $= \int_0^2e^{u}\,du = e^{2} - 1$.
\item $\int_0^1x\cdot x\,dx = \int_0^1x\,dx\cdot\int_0^1x\,dx = \frac14$.
\end{enumerate}
\end{karte}

\stupser{gruen}{LIATE: Logarithmus, Inverse, Algebraisch, Trigonometrisch, Exponentiell --- was weiter vorn steht, wird $u$. Bei (c) zweimal partiell, dann nach dem Integral auflösen.}{22pt}
\begin{karte}{gruen}{3.38}{\typwerk}{\meta{1}{15 min}{Band 3, Kap.\ 10, Aufgabe 10.1}}
\kartentitel{Partielle Integration}
\worum{Drei Klassiker: einmal, einmal mit $v' = 1$, einmal im Kreis.}
\methoden{Partielle Integration, Integral im Kreis}
Berechnen Sie (a) $\int xe^{x}\,dx$, (b) $\int\ln x\,dx$, (c) $\int_0^{\pi}e^{x}\sin x\,dx$.
\end{karte}

\stupser{gruen}{(a) $u = x^{2}$, Grenzen $0\to0$, $2\to4$. (b) $x = \sin t$. (c) $u = \ln x$.}{22pt}
\begin{karte}{gruen}{3.39}{\typwerk}{\meta{2}{15 min}{Band 3, Kap.\ 10, Aufgabe 10.2}}
\kartentitel{Substitution}
\worum{Drei Substitutionen: eine naheliegende, eine trigonometrische, eine, bei der man das Innere erst sehen muss.}
\methoden{Substitution (Integral)}
Berechnen Sie (a) $\int_0^2xe^{x^{2}}\,dx$, (b) $\int_0^1\sqrt{1 - x^{2}}\,dx$, (c) $\int\frac{dx}{x\ln x}$.
\end{karte}

\stupser{gruen}{$x^{2} - x - 2 = (x - 2)(x + 1)$. Ansatz $\frac A{x - 2} + \frac B{x + 1}$; $A$ durch Einsetzen von $x = 2$, $B$ mit $x = -1$.}{22pt}
\begin{karte}{gruen}{3.40}{\typwerk}{\meta{1}{10 min}{Band 3, Kap.\ 10, Aufgabe 10.3}}
\kartentitel{Partialbruchzerlegung}
\worum{Rationale Funktionen integriert man, indem man sie in Brüche mit linearem Nenner zerlegt. Die Zuhaltemethode macht das schnell.}
\methoden{Partialbruchzerlegung, Zuhaltemethode}
Berechnen Sie $\displaystyle\int\frac{3x + 5}{x^{2} - x - 2}\,dx$.
\end{karte}

\stupser{gruen}{Bogenlänge $\int\sqrt{1 + f'(x)^{2}}\,dx$ mit $f(x) = \sqrt{1 - x^{2}}$; der Integrand wird $\frac1{\sqrt{1 - x^{2}}}$ --- die Ableitung von $\arcsin$.}{22pt}
\begin{karte}{gruen}{3.41}{\typwerk}{\meta{2}{15 min}{Band 3, Kap.\ 10, Aufgabe 10.4}}
\kartentitel{Der Umfang des Einheitskreises}
\worum{$\pi$ wurde in Kapitel 5 über die Nullstellen des Kosinus definiert. Hier schließt sich der Kreis: Der Umfang ist tatsächlich $2\pi$.}
\methoden{Bogenlänge, Uneigentliches Integral}
Bestimmen Sie den Umfang des Einheitskreises über die Bogenlängenformel. Woran muss man am Rand des Intervalls denken?
\end{karte}

%%%%%%%%%%% Kapitel 11
\stupser{gruen}{$\int_1^R x^{-p}\,dx = \frac{R^{1 - p} - 1}{1 - p}$ für $p\neq1$. Was passiert für $R\to\infty$, je nach Vorzeichen von $1 - p$? Analog bei $0$.}{22pt}
\begin{karte}{gruen}{3.42}{\typwerk}{\meta{2}{10 min}{Band 3, Kap.\ 11, Aufgabe 11.1}}
\kartentitel{$\int x^{-p}$, zweimal uneigentlich}
\worum{Die Vergleichsintegrale für alles Weitere: bei $\infty$ konvergiert $x^{-p}$ für $p > 1$, bei $0$ für $p < 1$. Genau umgekehrt, und $p = 1$ fällt beide Male durch.}
\methoden{Uneigentliches Integral, Fallunterscheidung}
Untersuchen Sie $\int_1^{\infty}\frac{dx}{x^{p}}$ und $\int_0^1\frac{dx}{x^{p}}$ auf Konvergenz in Abhängigkeit von $p > 0$.
\end{karte}

\stupser{gruen}{(a) Partiell mit $u = \frac1x$, $v' = \sin x$: Randterm beschränkt, Restintegral $\int\frac{\cos x}{x^{2}}$ absolut konvergent. (b) Auf $[k\pi, (k + 1)\pi]$ ist $\int|\sin x|\,dx = 2$ und $\frac1x\ge\frac1{(k + 1)\pi}$.}{22pt}
\begin{karte}{gruen}{3.43}{\typwerk}{\meta{3}{20 min}{Band 3, Kap.\ 11, Aufgabe 11.2}}
\kartentitel{$\int\frac{\sin x}x$}
\worum{Das Integral, das konvergiert, aber nicht absolut --- die harmonische Reihe in Integralform.}
\methoden{Uneigentliches Integral, Partielle Integration, Abschätzung nach unten}
Untersuchen Sie (a) $\int_0^{\infty}\frac{\sin x}x\,dx$ und (b) $\int_0^{\infty}\frac{|\sin x|}x\,dx$ auf Konvergenz.
\end{karte}

\stupser{tief}{$e^{-x^{2}}\le e^{-x}$ für $x\ge1$. $\frac{\ln x}{x^{2}}\le\frac{x^{1/2}}{x^{2}}$ für große $x$. Für (c), (d) rechnet man direkt.}{22pt}
\begin{karte}{tief}{3.44}{\typtypisch}{\meta{2}{15 min}{Band 3, Kap.\ 11}}
\kartentitel{Vergleichen statt rechnen}
\worum{Bei uneigentlichen Integralen will man oft nur wissen, ob sie konvergieren. Dann vergleicht man mit $x^{-p}$ oder $e^{-x}$ --- und spart sich die Stammfunktion, die es manchmal gar nicht gibt.}
\methoden{Vergleichskriterium, Uneigentliches Integral, Grenzübergang}
\begin{enumerate}[label=(\alph*),leftmargin=1.6em,itemsep=2pt]
\item Zeigen Sie, dass $\int_0^{\infty}e^{-x^{2}}\,dx$ konvergiert, ohne es auszurechnen.
\item Berechnen Sie $\int_0^{\infty}\frac{dx}{1 + x^{2}}$.
\item Berechnen Sie $\int_1^{\infty}\frac{\ln x}{x^{2}}\,dx$.
\item Berechnen Sie $\int_0^1\ln x\,dx$. Wo ist dieses Integral uneigentlich?
\end{enumerate}
\end{karte}

%%%%%%%%%%% Kapitel 12
\stupser{akzent}{$f_n(x) = x^{n}$ auf $[0, 1]$: Grenzfunktion ist $0$ für $x < 1$ und $1$ bei $x = 1$. Der Abstand $\sup|f_n - f|$ bleibt $1$.}{22pt}
\begin{karte}{akzent}{3.45}{\typstolper}{\meta{3}{20 min}{Band 3, Kap.\ 12}}
\kartentitel{Punktweise ist nicht gleichmäßig}
\worum{Drei Dinge, die man gern vertauscht --- Grenzwert und Stetigkeit, Grenzwert und Integral, Grenzwert und Ableitung --- und die man nur bei gleichmäßiger Konvergenz vertauschen darf.}
\methoden{Gleichmäßige Konvergenz, Gegenbeispiel, Vertauschung von Grenzprozessen}
\begin{enumerate}[label=(\alph*),leftmargin=1.6em,itemsep=2pt]
\item $f_n(x) = x^{n}$ auf $[0, 1]$. Bestimmen Sie die Grenzfunktion. Ist sie stetig? Ist die Konvergenz gleichmäßig?
\item $g_n(x) = nx(1 - x^{2})^{n}$ auf $[0, 1]$. Zeigen Sie $g_n\to0$ punktweise, aber $\int_0^1g_n\to\frac12$.
\item $h_n(x) = \frac{x^{n}}n$ auf $[0, 1]$. Zeigen Sie, dass $h_n\to0$ gleichmäßig, aber $h_n'(1)\not\to0$.
\item Welcher Satz rettet welche Vertauschung, und unter welcher Voraussetzung?
\end{enumerate}
\end{karte}

\stupser{gruen}{$\sum x^{k} = \frac1{1 - x}$ gliedweise ableiten (erlaubt im Inneren des Konvergenzkreises), dann mit $x$ multiplizieren.}{22pt}
\begin{karte}{gruen}{3.46}{\typwerk}{\meta{1}{10 min}{Band 3, Kap.\ 12, Aufgabe 12.1}}
\kartentitel{$\sum kx^{k}$}
\worum{Was gleichmäßige Konvergenz praktisch bringt: Potenzreihen darf man gliedweise ableiten, und aus der geometrischen Reihe werden neue.}
\methoden{Gliedweise Ableitung, Geometrische Reihe}
Bestimmen Sie eine geschlossene Form für $\sum_{k\ge1}kx^{k}$ ($|x| < 1$) und berechnen Sie $\sum_{k\ge1}\frac k{2^{k}}$.
\end{karte}

%%%%%%%%%%% Kapitel 13
\stupser{tief}{$f_x = 2xy + ye^{xy}$, $f_y = x^{2} + xe^{xy}$. Richtungsableitung: Gradient mal Einheitsvektor. Tangentialebene: $z = f(p) + \nabla f(p)\cdot(x - p)$.}{22pt}
\begin{karte}{tief}{3.47}{\typtypisch}{\meta{2}{15 min}{Band 3, Kap.\ 13}}
\kartentitel{Gradient, Richtung, Tangentialebene}
\worum{Die Grundrechnungen in zwei Variablen an einem Beispiel: partielle Ableitungen, Gradient, Richtungsableitung, lineare Näherung.}
\methoden{Gradient, Richtungsableitung, Lineare Näherung}
$f(x, y) = x^{2}y + e^{xy}$, Punkt $p = (1, 0)$.
\begin{enumerate}[label=(\alph*),leftmargin=1.6em,itemsep=2pt]
\item Berechnen Sie $\nabla f(p)$.
\item Berechnen Sie die Richtungsableitung in Richtung $(1, 1)$. In welcher Richtung wächst $f$ am stärksten, und wie stark?
\item Geben Sie die Tangentialebene in $p$ an und schätzen Sie damit $f(1{,}1;\ 0{,}05)$.
\end{enumerate}
\end{karte}

\stupser{gruen}{$\nabla f = 0$ gibt $y = x^{2}$ und $x = y^{2}$, also $x = x^{4}$. Hesse-Matrix $\begin{pmatrix}6x & -3\\ -3 & 6y\end{pmatrix}$.}{22pt}
\begin{karte}{gruen}{3.48}{\typwerk}{\meta{2}{15 min}{Band 3, Kap.\ 13, Aufgabe 13.1}}
\kartentitel{Kritische Punkte klassifizieren}
\worum{Gradient null, Hesse-Matrix prüfen: Minimum, Maximum oder Sattel. Hier je einmal Sattel und Minimum.}
\methoden{Kritische Punkte, Hesse-Matrix}
Bestimmen und klassifizieren Sie alle kritischen Punkte von $f(x, y) = x^{3} - 3xy + y^{3}$.
\end{karte}

\stupser{gruen}{$\nabla(xy) = \lambda\nabla(x + y - 10)$: $y = \lambda$, $x = \lambda$.}{22pt}
\begin{karte}{gruen}{3.49}{\typwerk}{\meta{1}{10 min}{Band 3, Kap.\ 13, Aufgabe 13.2}}
\kartentitel{Lagrange am Rechteck}
\worum{Die Lagrange-Methode an einem Beispiel, dessen Antwort man kennt --- um zu sehen, wie die Methode sie findet.}
\methoden{Lagrange-Multiplikatoren, Nebenbedingung einsetzen}
Bestimmen Sie das Rechteck maximalen Flächeninhalts mit Umfang $20$ mit Lagrange-Multiplikatoren. Vergleichen Sie mit dem Weg über das Einsetzen der Nebenbedingung.
\end{karte}

\stupser{gruen}{Kugelkoordinaten, Funktionaldeterminante $r^{2}\sin\vartheta$; das Dreifachintegral zerfällt in drei Faktoren.}{22pt}
\begin{karte}{gruen}{3.50}{\typwerk}{\meta{2}{10 min}{Band 3, Kap.\ 13, Aufgabe 13.3}}
\kartentitel{Das Volumen der Kugel}
\worum{Transformationssatz und Fubini in einer Rechnung: $\frac43\pi$, und man sieht, woher jeder Faktor kommt.}
\methoden{Kugelkoordinaten, Transformationssatz, Satz von Fubini}
Berechnen Sie das Volumen der Einheitskugel im $\mathbb R^{3}$.
\end{karte}
